\documentclass[11pt]{article}

\usepackage{amsmath,amssymb}
\usepackage{xurl}
\usepackage[hidelinks]{hyperref}
\setlength{\emergencystretch}{2em}

% Shared commands for notes in docs/paper-gaps/.

\DeclareUnicodeCharacter{2080}{\ifmmode{}_{0}\else\textsubscript{0}\fi}
\DeclareUnicodeCharacter{2081}{\ifmmode{}_{1}\else\textsubscript{1}\fi}
\DeclareUnicodeCharacter{2082}{\ifmmode{}_{2}\else\textsubscript{2}\fi}
\DeclareUnicodeCharacter{2099}{\ifmmode{}_{n}\else\textsubscript{n}\fi}

\newcommand{\Lean}{\textsc{Lean}}
\ExplSyntaxOn
\NewDocumentCommand{\leanid}{m}{
  \ifmmode
    \textsf{\detokenize{#1}}
  \else
    \group_begin:
    \urlstyle{sf}
    \tl_set:Nn \l_tmpa_tl {#1}
    \tl_replace_all:Nnn \l_tmpa_tl {\_} {_}
    \exp_args:NV \path \l_tmpa_tl
    \group_end:
  \fi
}
\ExplSyntaxOff
\providecommand{\tr}{\operatorname{tr}}
\newcommand{\tnleanIssueBase}{https://github.com/LionSR/TNLean/issues/}
\newcommand{\tnleanPrBase}{https://github.com/LionSR/TNLean/pull/}
\newcommand{\tnleanDocBase}{https://sirui-lu.com/TNLean/docs/}
\newcommand{\ghissue}[1]{\href{\tnleanIssueBase#1}{GitHub issue \##1}}
\newcommand{\ghpr}[1]{\href{\tnleanPrBase#1}{PR \##1}}
\newcommand{\doclink}[1]{\href{\tnleanDocBase#1.html}{\path{#1}}}

% Dirac notation and the identity operator, as in blueprint/src/macros/common.tex.
% \providecommand keeps note-local definitions authoritative.
\usepackage{dsfont}
\providecommand{\ket}[1]{|#1\rangle}
\providecommand{\bra}[1]{\langle #1|}
\providecommand{\braket}[2]{\langle #1 | #2 \rangle}
\providecommand{\ketbra}[2]{|#1\rangle\!\langle #2|}
\providecommand{\Id}{\mathds{1}}
\providecommand{\one}{\mathds{1}}
\providecommand{\C}{\mathbb{C}}
\providecommand{\MN}[1]{M_{#1}(\C)}
\providecommand{\MD}{M_D(\C)}

% External referencing for paper sources.  The build helper writes sanitized
% auxiliary files under build/paper-aux/ and excludes duplicated source labels.
\usepackage{xr}

% Import a generated paper auxiliary file with a caller-supplied label prefix.
% The candidate paths support builds from docs/paper-gaps/, the repository root,
% and build/paper-aux/ itself.
\newcommand{\importpaperaux}[2]{%
  \IfFileExists{../../build/paper-aux/#2.aux}{%
    \externaldocument[#1]{../../build/paper-aux/#2}%
  }{%
    \IfFileExists{build/paper-aux/#2.aux}{%
      \externaldocument[#1]{build/paper-aux/#2}%
    }{%
      \IfFileExists{#2.aux}{%
        \externaldocument[#1]{#2}%
      }{}%
    }%
  }%
}

% General external reference macros
\newcommand{\paperref}[2]{\ref{#1-#2}}
\newcommand{\papereqref}[2]{(\ref{#1-#2})}

% CPSV16 source (1606.00608 / MPDO-22-12-17-2.tex) external document and shortcuts
\importpaperaux{cpsv16-}{1606.00608/MPDO-22-12-17-2-tnlean}
\newcommand{\cpsvref}[1]{\paperref{cpsv16}{#1}}
\newcommand{\cpsveqref}[1]{\papereqref{cpsv16}{#1}}

% Verdict marker: \gapnote{<kind>}{<status>}. Kind is the policy
% classification (clarification, local-correction, scope-restriction,
% unfaithful, false-source, open-gap); status is open, wip (elimination
% underway), resolved, or historical. Machine-read by the site index;
% prints a verdict line.
\newcommand{\gapnote}[2]{\par\noindent\textbf{Verdict:} #1 (#2).\par}


\title{Theorem Surface Audit for Canonical Form, BNT, and the Non-periodic MPS Fundamental Theorem}
\author{}
\date{2026-05-05}

\begin{document}
\maketitle

\section{Purpose}

This note compares the theorem surface in the formal development with the
statements in the source papers.  It is not an issue index.  The goal is to
separate four kinds of results:
\begin{enumerate}
  \item named statements appearing directly in the papers;
  \item cited mathematical results used by the papers;
  \item auxiliary formal statements needed to make an implicit proof step exact;
  \item stronger or alternative statements that should not be presented as part
        of the paper-level proof unless they become necessary.
\end{enumerate}

\section{Named source statements}

The 2016 source has a small number of named statements in the canonical-form
and BNT part:
\begin{itemize}
  \item normal tensor and canonical form definitions
    \cite[lines~233--255]{Cirac2016MPDO_arXiv};
  \item basis of normal tensors definition and characterization
    \cite[lines~271--280]{Cirac2016MPDO_arXiv};
  \item injective and block-injective canonical form definitions, together with
    the quantum-Wielandt and block-injective-after-blocking assertions
    \cite[lines~317--345]{Cirac2016MPDO_arXiv};
  \item the proportional Fundamental Theorem and the equal-MPV corollary
    \cite[lines~347--356]{Cirac2016MPDO_arXiv}.
\end{itemize}
The 2021 review gives the same normal tensor, BNT, proportional theorem, and
equal-MPV corollary in Section~IV
\cite[lines~1828--1899]{Cirac2021Matrix}.

These are the paper-level targets.  They should remain the prominent theorem
statements in the blueprint.

\section{Why the formalization has more declarations}

The paper proof suppresses finite-dimensional and finite-index operations that
must be named in a proof assistant.  In this part of the development the extra
formal statements usually fall into one of the following categories.

\paragraph{Reduction statements.}
The canonical-form proof splits the paper's ``put the tensor in canonical
form'' step into invariant-subspace extraction, block-diagonal decomposition,
trace-preserving normalization, period-removing blocking, and primitive-block
normality.  These are not new mathematical claims beyond the source reduction,
but they need separate hypotheses and conclusions so later theorems can apply
them.

\paragraph{Common blocking statements.}
The paper says that periodicity is removed by blocking a common multiple of the
periods.  The formalization must record positivity of the chosen block length,
persistence of injectivity under further blocking, and the simultaneous
application to a finite family of representative sectors.  These statements are
auxiliary, but they are mathematically faithful to the period-removal step.

\paragraph{Wielandt proof-chain statements.}
The source proof of the quantum-Wielandt bound uses Lemma~1, eigenvalue
extraction, vector spreading, and the fixed-length rank-one step as an argument,
not as separate MPS theorem statements
\cite[Theorem~1 proof]{Sanz2010Quantum}.  Formal conjunctions that bundle
these inputs, such as the \emph{Wielandt analysis} and \emph{Wielandt chain}
declarations, should therefore be lemmas.  The visible theorem is the
cumulative or fixed-length Wielandt conclusion, with the stronger fixed-length
claim distinguished from the cumulative span bound.

\paragraph{BNT comparison statements.}
The source BNT characterization selects one representative from each
phase-gauge equivalence class and compares the power sums of weights.  The
formalization separates this into overlap decay, eventual linear independence,
permutation rigidity, sector-weight comparison, and phase absorption.  The
heterogeneous block-matching theorem is a genuine formal endpoint for the
block-matching part, while the final global weighted gauge equivalence remains
the paper-level equal-MPV target.
The overlap-to-linear-independence and change-of-basis declarations are
linear-algebra support lemmas for the permutation-rigidity theorem, not
additional BNT theorems from the source.

\section{External lemmas need explanations}

When a result is needed but the MPS source only cites it, the formal
development should give a pedagogical explanation at the point where the result
enters the proof.  A citation alone is not enough in this situation.  The
explanation should include:
\begin{enumerate}
  \item the exact mathematical statement used;
  \item the hypotheses in the MPS notation, after blocking or reindexing if
        necessary;
  \item a short proof idea, or a precise indication that the proof is imported
        from the cited source;
  \item the formal declaration that supplies the result, if it has already been
        formalized.
\end{enumerate}

This applies in particular to the following families of external input.
\begin{itemize}
  \item The quantum Wielandt theorem supplies finite-length matrix-algebra span
    from normality or primitivity \cite{Sanz2010Quantum}.  The paper cites it in
    the canonical-form reduction; the formal text should say whether it is using
    cumulative span, exact-length span, injectivity of the map $\Gamma_L$, or
    block-injectivity of a direct sum.
  \item The block-injective canonical-form bound uses the older MPS
    injectivity and block-injectivity argument of
    \cite{PerezGarcia2007MPSReps}.  The formal statement should identify the
    common blocked physical alphabet and the direct-sum matrix algebra that is
    being spanned.  The direct-sum proof input in the older source is recorded
    in \path{docs/paper-gaps/david2006_direct_sum_input.tex}.
  \item Vandermonde inversion and Newton--Girard identities are elementary, but
    they are not proved in the MPS source.  When they are used to recover
    coefficients, multiplicities, or weight multisets, the blueprint should
    explain the finite linear system or the power-sum recursion.
  \item Perron--Frobenius and peripheral-spectrum facts for quantum channels are
    standard inputs from finite-dimensional channel theory, but they are not
    interchangeable.  A formal theorem should specify whether it uses
    irreducibility, primitivity, trace preservation, a faithful fixed point, or
    a rank-one peripheral projection.
\end{itemize}

This rule also gives a retirement test.  If an external-looking lemma has no
source citation, no proof explanation, and no current proof path needing it, it
should not remain as a headline theorem in the paper-level blueprint.

\section{Current demotions}

The following declarations are auxiliary and should not be read as independent
paper-level theorems.
\begin{itemize}
  \item The explicit-coefficient equal-MPV result is named
    \[
      \texttt{fundamentalTheorem\_equalMPV\_of\_explicit\_coefficients}.
    \]
    It is not the full equal-MPV corollary from the papers.  It assumes the
    coefficient arrays and their nonzero limits.
  \item The proportional result
    \[
      \texttt{fundamentalTheorem\_proportionalMPV\_CFBNT}
    \]
    is stated with the same explicit coefficient arrays, their limits, and
    non-vanishing assumptions.  The source theorem derives the relevant BNT
    coefficient comparison from the proportional MPV hypothesis; that extraction
    is not part of this formal statement.
  \item The unused fixed-size BNT Vandermonde MPV declarations have been
    removed.  The scalar Vandermonde separation lemma remains because it is a
    small algebraic input elsewhere, but it is not presented as part of the
    main BNT comparison route.
  \item The equal-span permutation rigidity statement with asymptotically
    orthonormal overlaps is a lemma.  It is a useful span route, but the
    proportional BNT permutation theorem is the theorem used by the main
    proportional Fundamental Theorem in Chapter~11.
  \item The older theorem from proportional-decomposition data to sector
    comparison without leftover-zero-block bookkeeping has been removed.  The
    active after-blocking route keeps this bookkeeping visible because the
    formal equality relation includes the empty word.  This is not terminology
    from \cite{Cirac2016MPDO_arXiv}; the source only says that canonical forms
    may have zero blocks through the inequality
    \(\sum_k D_k \le D\) \cite[lines~217--219]{Cirac2016MPDO_arXiv}.
\end{itemize}

\section{Statements to keep visibly conditional}

Several declarations should stay in the development, but their statements
should remain visibly conditional until the missing paper input is produced.
\begin{itemize}
  \item The conditional after-blocking Fundamental Theorem assumes the
    post-blocking sector hypotheses.  It should not be cited as the unconditional
    canonical-form theorem.
  \item The explicit-coefficient equal case assumes coefficient arrays rather
    than deriving the BNT power-sum data from bare equality of MPVs.
\end{itemize}

\section{Retirement criterion}

An auxiliary declaration should be removed or moved out of the main blueprint
route when all three conditions hold:
\begin{enumerate}
  \item it is not a statement in the cited paper and is not a cited outside
        result;
  \item no current proof of a paper-level target uses it;
  \item it states a stronger finite-length, span, or synchronization hypothesis
        than the source proof requires.
\end{enumerate}
If a declaration is still used but only as a formal intermediate result, it
should be a lemma or a clearly conditional theorem, not a headline replacement
for the paper theorem.

\bibliographystyle{alpha}
\bibliography{docs/paper-gaps/references}

\end{document}
